The KKT slack from \cref{sec:obstacle-active-sets} has an exact network
interpretation. Form an augmented graph by
adding a ground vertex $\mathsf g$.  Give every original edge conductance
$(1-\alpha)/2$ and connect each $i\in\gV$ to $\mathsf g$ with conductance
$\alpha d_i$; when $\alpha=1$, omit the zero-conductance original edges.  Fix
an orientation, let $\mB$ be the signed incidence matrix with the ground
column deleted, and let $\mW$ be the diagonal matrix of positive
conductances.  Define
\begin{equation}
 \mH:=\mD^{1/2}\mQ\mD^{1/2}
 =\alpha\mD+\frac{1-\alpha}{2}(\mD-\mA),
 \qquad
 \vc_\rho^{\mathrm d}:=\mD^{1/2}\vc_\rho=\alpha(\vs-\rho\mD\one).
 \label{eq:grounded-system}
\end{equation}
Then $\mB^\trans\mW\mB=\mH$.

\begin{proposition}[Grounded electrical-flow dual]
\label{prop:grounded-flow-dual}
Under the change of variables $\vy=\mD^{-1/2}\vx$, the obstacle problem is
\begin{equation}
 \min_{\vy\geq\vzero}
 \left\{
  \frac12\norm{\mW^{1/2}\mB\vy}_2^2-(\vc_\rho^{\mathrm d})^\trans\vy
 \right\}.
 \label{eq:grounded-flow-primal}
\end{equation}
Its Fenchel dual, written as a minimum-energy problem, is
\begin{equation}
 \min_{\vf}\ \frac12\vf^\trans\mW^{-1}\vf
 \qquad\text{subject to}\qquad
 \mB^\trans\vf\geq\vc_\rho^{\mathrm d}.
 \label{eq:grounded-flow-dual}
\end{equation}
The two minimum values have opposite signs.  At an optimal primal--dual pair,
\begin{equation*}
 \vf^*=\mW\mB\vy^*,
 \qquad
 \vlambda^*:=\mB^\trans\vf^*-\vc_\rho^{\mathrm d}
 =\mD^{1/2}\vw^*\geq\vzero,
 \qquad
 y_i^*\lambda_i^*=0.
\end{equation*}
More generally, let $\vx^\gU$ be the zero-extended exact solution on a
reachable active set $\gU$, with slack $\vw^\gU$ as defined in
\cref{subsec:reachable-faces}. Set $\vy^\gU:=\mD^{-1/2}\vx^\gU$.
The induced flow $\vf^\gU:=\mW\mB\vy^\gU$ is tight on $\gU$ and
\begin{equation*}
 (\mB^\trans\vf^\gU-\vc_\rho^{\mathrm d})_j=\sqrt{d_j}\,w_j^\gU
 \quad(j\notin \gU).
\end{equation*}
Consequently, $w_j^\gU<0$ is exactly a violated dual vertex constraint.
\end{proposition}

\begin{proof}
The identities in \cref{eq:grounded-system} give
$\vy^\trans\mH\vy=\norm{\mW^{1/2}\mB\vy}_2^2$ and transform
\cref{eq:obstacle-problem} into \cref{eq:grounded-flow-primal}.  The conjugate
identity
\[
 \frac12\norm{\mW^{1/2}\vz}_2^2
 =\max_{\vf}
 \left\{\vf^\trans\vz-\frac12\vf^\trans\mW^{-1}\vf\right\}
\]
and minimization over $\vy\geq\vzero$ make the inner value finite precisely
when $\mB^\trans\vf-\vc_\rho^{\mathrm d}\geq\vzero$.  Fenchel--Rockafellar duality applies
because the quadratic is continuous everywhere and the orthant is nonempty,
with the sign reversal in the displayed minimization convention.  Stationarity in the conjugate pair
gives $\vf^*=\mW\mB\vy^*$; minimization over the orthant gives
complementarity.  Finally,
\[
 \mB^\trans\mW\mB\vy-\vc_\rho^{\mathrm d}
 =\mH\vy-\vc_\rho^{\mathrm d}
 =\mD^{1/2}(\mQ\vx-\vc_\rho)
 =\mD^{1/2}\vw,
\]
which proves both slack identities.
\end{proof}

Thus the active-set method can be interpreted as local constraint
generation for \cref{eq:grounded-flow-dual}: a restricted solve constructs
the flow induced by the current potentials, and boundary slack tests identify
violated vertex inequalities. The implementation uses the local support
and certificate properties of \cref{sec:obstacle-active-sets} to perform
these tests; the batch-depth bound in \cref{sec:threshold-batch-depth}
controls the number of restricted solves. Constructing the full dual or
testing every vertex would require graph-wide access.
